% 21-testing.tex — Chapter 21: Testing Infrastructure
\chapter{Testing Infrastructure}
\label{chap:21-testing}
\index{testing infrastructure}
\index{test suite}

\pnum
The Lain compiler's correctness is verified by a test suite of Lain source files
under \srcref{tests/}.  Tests are driven by the shell script
\srcref{run\_tests.sh}.

% -------------------------------------------------------------------------
\section{Test file conventions}
\label{sec:test-conventions}
\index{test file conventions}
% -------------------------------------------------------------------------

\pnum
Test files follow a naming convention that determines the expected outcome:

\begin{center}
\begin{tabular}{ll}
\toprule
\textbf{Name pattern} & \textbf{Expected outcome} \\
\midrule
\texttt{*\_pass.ln} & Compilation must succeed (exit code 0). \\
\texttt{*\_fail.ln} & Compilation must fail (exit code non-zero). \\
Other \texttt{.ln}  & Treated as pass tests by default. \\
\bottomrule
\end{tabular}
\end{center}

\pnum
A \texttt{\_fail.ln} test may additionally contain a comment of the form
\texttt{// EXPECT: [EXXX]} in the source text.  If present, the test framework
verifies that the compiler's \texttt{stderr} output contains the expected
diagnostic code.  This allows negative tests to verify not just that the
compiler rejects the program but that it emits the \emph{correct} error.

% -------------------------------------------------------------------------
\section{Test runner}
\label{sec:test-runner}
% -------------------------------------------------------------------------

\pnum
\cname{run\_tests.sh} (\srcref{run\_tests.sh}):
\begin{enumerate}
\item Builds the compiler if not already present.
\item Recursively discovers all \texttt{.ln} files under \srcref{tests/}.
\item For each file, runs the compiler and compares the exit code against the
  expected outcome.
\item For \texttt{\_fail.ln} files with \texttt{// EXPECT:} comments, also
  matches the diagnostic code against \texttt{stderr}.
\item Reports pass/fail counts and lists all failed tests.
\end{enumerate}

% -------------------------------------------------------------------------
\section{Test categories}
\label{sec:test-categories}
\index{test categories}
% -------------------------------------------------------------------------

\pnum
Tests are organized into subdirectories under \srcref{tests/}:

\begin{center}
\begin{tabular}{ll}
\toprule
\textbf{Directory} & \textbf{Tests} \\
\midrule
\texttt{tests/safety/}     & Ownership, linearity, borrow, and unsafe tests \\
\texttt{tests/safety/bounds/} & Bounds checking (pass and fail) \\
\texttt{tests/safety/purity/} & Purity / termination verification \\
\texttt{tests/types/}      & Type inference, widening, mismatch detection \\
\texttt{tests/emit/}       & Code generation correctness \\
\texttt{tests/stdlib/}     & Standard library function tests \\
\texttt{tests/generics.ln} & Comptime generic instantiation \\
\texttt{tests/core/}       & Core language features \\
\texttt{tests/borrowck/}   & Borrow checking pass and fail tests \\
\bottomrule
\end{tabular}
\end{center}

% -------------------------------------------------------------------------
\section{Invoking the test suite}
\label{sec:test-invoke}
% -------------------------------------------------------------------------

\pnum
To run the full test suite from the repository root:

\begin{ccode}
bash run_tests.sh
\end{ccode}

\pnum
The script builds the compiler automatically if the \texttt{lain} binary is not
present.  Exit code 0 means all tests passed; non-zero means at least one test
failed.  Failed test names are printed to \texttt{stdout}.

\pnum
To run a single test manually:

\begin{ccode}
./lain tests/safety/bounds/dynamic_len_known_idx_fail.ln
echo "exit: $?"
\end{ccode}

\pnum
The expected exit code for a \texttt{\_fail} test is 1 (or any non-zero value).
For a \texttt{\_pass} test, exit code 0 is required.
